\section*{Aktivitas Pembelajaran dan Praktikum}
\addcontentsline{toc}{section}{Aktivitas Pembelajaran dan Praktikum}

\begin{table}[H]
\centering
\caption{Rencana aktivitas Bab XI}
\label{tab:ch11-aktivitas}
\begin{tabularx}{\textwidth}{@{}llX@{}}
\toprule
\textbf{Sesi} & \textbf{Bentuk} & \textbf{Kegiatan} \\
\midrule
1 & Ceramah & Model shift-reduce dan handle \\
2 & Latihan & Menelusuri stack untuk ekspresi \\
3 & Lokakarya & Membangun tabel operator precedence \\
4 & Praktikum & Kalkulator operator precedence \\
\bottomrule
\end{tabularx}
\end{table}

\subsection*{Lokakarya: Penelusuran dan Kalkulator}
\textbf{Tujuan.} Memahami shift-reduce melalui penelusuran manual
dan implementasi kalkulator sederhana.

\textbf{Langkah kerja.}
\begin{enumerate}[label=\arabic*.,leftmargin=*]
  \item Telusuri shift-reduce untuk \(\code{id} * \code{id} +
        \code{id}\).
  \item Tandai handle pada setiap langkah reduce.
  \item Bangun tabel operator precedence untuk operator
        \(+\), \(-\), \(*\), \(/\), dan tanda kurung.
  \item Modifikasi \file{shift\_reduce.c} agar mendukung operator
        \(-\), \(/\), dan tanda kurung.
  \item Uji dengan ekspresi yang mengandung precedence dan
        associativity berbeda.
\end{enumerate}

\textbf{Definisi selesai.}
\begin{checklist}
  \item Penelusuran manual benar.
  \item Tabel precedence lengkap.
  \item Kalkulator mendukung empat operator dan kurung.
  \item Minimal 10 ekspresi uji.
  \item Data penelusuran dan tabel precedence siap dipakai pada
        Bab XII.
\end{checklist}

Bab ini tidak menambahkan milestone proyek baru. Hasil lokakarya
menjadi data pendukung Praktikum 3 (12\%) yang dilaporkan pada
Bab XIII. Tabel ACTION/GOTO SLR/LALR dibangun pada Bab XII.
